\chapter{Практические рекомендации для компании}

\section{Управленческий cash flow}

Компании необходимо перейти от разового анализа к ежемесячному управленческому cash flow. Классификация банковских операций должна выполняться регулярно, а внутренние переводы, кредиты и операционные платежи должны разделяться в отчетности.

\section{Контроль закупок и постоянных расходов}

Производственные закупки являются главной зоной оптимизации. Их нужно анализировать по поставщикам, ценам, объемам и условиям оплаты. Постоянные расходы, такие как аренда и коммунальные платежи, должны быть отделены от переменных закупок, чтобы видеть базовый cash burn.

\section{Платежный календарь}

Для месяцев риска необходимо сформировать платежный календарь. В него следует включить аренду, налоги, зарплаты, кредиты, регулярные сервисы и крупные закупки. Такой календарь также может быть внешним признаком для модели cash flow.

\section{Нормализация контрагентов}

Для надежной сверки банка и 1С требуется нормализатор контрагентов. Он должен учитывать разные юридические лица, расчетные счета, адреса торговых точек и округления сумм. Без этого автоматическая сверка будет давать ложные расхождения.

\section{Расширение клиентской базы}

Высокая зависимость от одного канала продаж создает риск для денежного потока. Компании следует расширять клиентскую базу и отслеживать концентрацию оплат не только по юридическим лицам, но и по брендам или сетям.

\section{Внедрение dashboard}

Результаты исследования целесообразно вывести в dashboard: cash flow по месяцам, прогноз остатка на 30 дней, операции на проверку, концентрация поставщиков и покупателей, прогноз спроса и возвратов, рекомендации по производству.
